% Compact core statement; full derivation is in integral_value_proof.tex.
\subsection{An integral correction with a complete price certificate}
\label{sec:integral-core-certificate}

I retain the classical risk-neutral basket model and alter the representation
of its pricing correction.  Let
\[
 P(m)=H\int_0^1\Phi\!\left(\frac{m-\log(1+Hs/K)}{\lambda}\right)ds,
 \quad T(y)=\sum_iw_ie^{d_i+s_i\cdot y},\quad U=T-1,\quad m=m_0+a\cdot y,
\]
where $w_i>0$, $\sum_iw_i=1$, $\sum_iw_id_i=0$ and $\sum_iw_is_i=0$.
Jensen's inequality gives $U\geq0$.  The undiscounted price is
$V=c_0+c_1+\mathbb E R_2$, with
$R_2=P(m+\log T)-P(m)-UP'(m)$ and the same continuous coefficients
$c_0,c_1$ used by the direct second-correction comparator.  Define
\[
 A_2(x)=-\lambda^{-1}-x\lambda^{-2},\quad
 D_2\geq\sup_x|A_2(x)\phi(x)|,\quad
 B_2=HD_2U_*^2/2,
\]
where $U_*>0$ bounds $U$ on the actual retained financial nodes.  I obtain
$U_*$ by outward evaluation of the $2^p$ vertices of their covering box:
coordinatewise convexity of $T$ proves that its maximum is at a vertex.
This finite normaliser does not replace the continuous-tail certificate.

Differentiation of $g_m(t)=P(m+\log t)$ and the integral Taylor remainder give
\begin{equation}
 R_2(y)=\frac{HU(y)^2}{2}\mathbb E_{Z,S}
 \left[\frac{A_2(\xi)\phi(\xi)}{(1+ZU(y))^2}\right],\qquad
 \xi=\frac{m+\log(1+ZU)-\log(1+HS/K)}{\lambda},
 \label{eq:core-integral-correction}
\end{equation}
where $S$ is uniform and $Z$ is the minimum of two independent uniforms.
Thus $Z$ has density $2(1-z)$.  With
\[
 a=(U/U_*)^2(1+ZU)^{-2},\qquad
 h(x)=\{1+A_2(x)\phi(x)/D_2\}/2,
\]
both $a,h$ lie in $[0,1]$ and
$p=(1-a)/2+ah(\xi)$ is a valid probability.
Its signed score is $X=2p-1$, with
$|X|\leq U^2/U_*^2$.  The matched direct comparator computes
$p^{\mathrm{dir}}=\{1+R_2/B_2\}/2$ with the same normaliser; it satisfies
the same pointwise score bound.  A further comparator uses the nonnegative
first correction $R_1=P(m+\log T)-P(m)$, bounded by
$B_1=\min\{H,HD_1U_*\}$.

I prepare independent uniform labels by Hadamards and compute actual dyadic
midpoint truncated-normal coordinates by reversible lookup.  All financial
scalar operations are enclosed endpoint tables with unsigned magnitude and
an explicit decreasing bit.  For a $C^2$ scalar on an interval of length $s$,
$2^b$ equal interpolation cells, endpoint error $e_*$ and output unit $2^{-q}$,
\[
 e_f\leq e_*+\|f''\|_\infty s^2/(8\,2^{2b})+2^{-q}.
\]
The proof is the two-endpoint interpolation remainder plus convex propagation
of nodal errors and one integer truncation unit.  Argument perturbation adds
$\|f'\|_\infty e_x$.  Products use absolute range-dependent errors, for
example $|\widehat x\widehat y-xy|\leq |x|e_y+|\widehat y|e_x$ before
truncation.  Every actual scalar domain, integer range and signed final
subtraction is checked.  Final projection onto $[0,1]$ is nonexpansive.
These checks yield a probability error $\eta_p$ and monetary arithmetic error
$2B_2\eta_p$, with no hidden tail or grid allowance.

Two exact dyadic transformations,
$p_+=\max(0,2p-1)$ and $p_-=\max(0,1-2p)$, satisfy $p_+-p_-=2p-1$.
If $M_{2,N}$ bounds the exact finite-law moment $\mathbb E U^2$, each
implemented flag has the prior
\begin{equation}
 \mathbb E\widehat p_\pm\leq
 \min\{1,M_{2,N}/U_*^2+2\eta_p\}.
 \label{eq:core-rare-prior}
\end{equation}
Indeed the ideal flag is at most $|X|$ and the positive-part map is
one-Lipschitz.  The flag difference remains exact, so its combined arithmetic
error is $2B_2\eta_p$.  I compute $M_{2,N}$ by the multinomial expansion of
$T^2$ and products of one-dimensional finite exponential sums, without
requiring a financial price.  Both flags and their joint statistical
confidence are charged.  This prior alone is not a physical-cost advantage.
The positive first-correction comparator receives its own finite mean prior,
using the smallest of $H\mathbb E\omega$, $HD_1\mathbb E[\omega U]$ and
$K e^{m_0+\lambda^2/2}\mathbb E[\omega e^{a\cdot Y}U]$, divided by $B_1$
and enlarged by its arithmetic error.  The last expectation is evaluated by
the same product sums and retains the level-ratio correlation; $\omega=1$
for the quantile construction.

Let $M_{3,N}$ bound the third finite moment, $D_3$ the third joint norm and
$E_2\geq\sup|(A_2\phi)'|$.  Coupling auxiliaries to cell midpoints gives
\[
 \beta_Z\leq HD_3M_{3,N}/2^{b_Z+2},\qquad
 \beta_S\leq H^2E_2M_{2,N}/(K\lambda\,2^{b_S+2}).
\]
Let $\beta_Y$ certify the original continuous correction against the
mass-weighted finite financial correction, $e_c$ bound the coefficient error,
and $[F_-,F_+]$ enclose the external Gaussian mass factor, whose midpoint is
used.  With $e_F=(F_+-F_-)/2$, the full deterministic undiscounted error is
\begin{equation}
 e_{\mathrm{det}}\leq e_c+\beta_Y+
 F_+(\beta_Z+\beta_S+2B_2\eta_p)+e_FB_2.
 \label{eq:core-monetary-certificate}
\end{equation}
This follows by inserting the exact coefficient, finite financial correction,
finite auxiliary average and implemented bounded score, and applying the
triangle inequality.  The direct second correction omits the auxiliary
terms; the first correction also replaces $2B_2\eta_p$ by $B_1\eta_p$.
Statistical error is added separately, and discounting multiplies the whole
bound to obtain the price $\Pi$.  The retained one-dimensional financial
transfer uses an already validated continuous reference and charges that
classical premise explicitly.  Higher-dimensional grid-only envelopes are
not presented as complete pricing certificates.  The full derivation, node
records and rejected budgets are retained in the supplementary material.
